\documentclass[11pt,a4paper]{article}
\usepackage[T1]{fontenc}
\usepackage{lmodern,amsmath,amssymb,amsthm,mathtools,booktabs,microtype}
\usepackage[margin=27mm]{geometry}
\usepackage[colorlinks=true,linkcolor=blue,citecolor=blue,urlcolor=blue]{hyperref}
\setlength{\emergencystretch}{2em}
\newtheorem{theorem}{Theorem}[section]
\newtheorem{lemma}[theorem]{Lemma}
\newtheorem{proposition}[theorem]{Proposition}
\newtheorem{corollary}[theorem]{Corollary}
\theoremstyle{definition}\newtheorem{definition}[theorem]{Definition}
\newcommand{\Q}{\mathbb Q}\newcommand{\F}{\mathbb F}
\DeclareMathOperator{\Nm}{N}
\title{Local solubility at every shared pole parameter}
\author{}\date{}
\begin{document}\maketitle
\begin{abstract}
Every rational value of the normalized conic parameter gives an everywhere
locally soluble second-pole quartic for every prepared input. The assertion
holds for an arbitrary finite family at one prescribed parameter, and local
witnesses can avoid the nonzero guards used by the finite-algebra encoding.
We give explicit proofs at all primes, including the primes dividing the
conic denominator, and an exact elliptic formulation of the remaining global
existence problem. No local obstruction can exclude the common-parameter
construction on prepared inputs. Simultaneous rational realization is not
asserted.
\end{abstract}

\section{The local theorem}

\begin{definition}
A rational number $b$ is \emph{prepared} if its reduced denominator is a
square and
\[
 0<b<17,\qquad b\in\Q_2^{\times2}\cap\Q_3^{\times2}
 \cap\Q_{17}^{\times2},\qquad
 v_2(b)\ge3,\quad v_3(b)\ge1,\quad v_{17}(b)\ge2.
\]
These are the preparation conditions of the pole source in
\cite{Fixed,UnitNorm}.
\end{definition}

Write $b=dt^2$ with $d,t\in\Q^\times$. The factorization may be the
syntactic square extraction used by the source; no squarefreeness is needed.
Set
\begin{equation}\label{eq:data}
 D=17d,\qquad C=17(17-b),\qquad
 e_\tau=\frac{\tau^2-D}{\tau^2+3D},\qquad A_\tau=e_\tau C.
\end{equation}
The normalized second-pole equation is
\begin{equation}\label{eq:quartic}
 \mathcal C_\tau:\qquad
 q^2=D(s^2-A_\tau)(s^2+3A_\tau).
\end{equation}
Its smooth projective model is understood when discussing points at infinity.

\begin{theorem}\label{thm:local}
For a prepared $b$ and every $\tau\in\Q$, the quantities in
\eqref{eq:data} are defined, $D A_\tau\ne0$, and
$\mathcal C_\tau$ is a smooth genus-one curve. It has points over every
completion of $\Q$. At each completion it has infinitely many affine points
with $q\ne0$, and these can avoid the zeros of any prescribed finite
collection of nonzero polynomials in $s$.
\end{theorem}

The parameter is arbitrary. Thus the same rational $\tau$ can be used
locally for every member of any finite family of prepared inputs, even when
their input square classes are distinct.

\begin{lemma}\label{lem:open}
Let $k$ be $\mathbb R$ or a finite extension of $\Q_p$, and let $D,A\in
k^\times$. If the smooth projective curve
$q^2=D(s^2-A)(s^2+3A)$ has a $k$-point, then it has infinitely many affine
$k$-points with $q\ne0$ outside any prescribed finite set of values of $s$.
\end{lemma}
\begin{proof}
The four branch points over an algebraic closure are distinct in
characteristic zero. Every projective $k$-point is therefore smooth, and a
local coordinate identifies a neighborhood with a nonempty open subset of
$k$, by the implicit function theorem. Such a neighborhood is infinite.
The points at infinity, the zeros of $q$, and the fibers over finitely many
values of $s$ form a finite set. Deleting them proves the assertion.
\end{proof}

\section{Proof at every place}

Since $b>0$, one has $d>0$ and $D>0$. Hence $\tau^2+3D>0$ for
$\tau\in\Q$. Preparation and $b=dt^2$ imply that $v_{17}(d)$ is even,
so $v_{17}(D)$ is odd. Thus $D$ is not a rational square and
$\tau^2\ne D$. Also $C>0$. Consequently $D A_\tau\ne0$.
The polynomial on the right of \eqref{eq:quartic} has discriminant
\begin{equation}\label{eq:disc}
                         -2^{12}3D^6A_\tau^6,
\end{equation}
which proves smoothness and genus one. Fix $\tau\in\Q$, and abbreviate
$e=e_\tau$, $A=A_\tau$.

\subsection{Primes other than $2,3,17$}

Fix such a prime $p$. If $D\in\Q_p^{\times2}$, the curve has two
$\Q_p$-points at infinity: in the chart $u=1/s$, $z=q/s^2$, the
equation at $u=0$ is $z^2=D$. If $-3D\in\Q_p^{\times2}$,
the value $s=0$ gives the nonzero square $-3DA^2$ on the right.

Suppose next that $v_p(D)$ is odd. The valuations of $D$ and $b$ have
the same parity. The reduced denominator of $b$ is a square, so a negative
valuation of $b$ is even. It follows that $v_p(b)>0$, and
\[
                      C=17^2(1-b/17)\in\Q_p^{\times2}.
\]
Write $D=h^2D_0$ with $h\in\Q_p^\times$ and $v_p(D_0)=1$, and
replace $\tau$ by $\tau/h$ in the formula for $e$. This leaves $e$
unchanged. If $v_p(\tau/h)\le0$, then
\begin{equation}\label{eq:units}
 e-1=\frac{-4D_0}{(\tau/h)^2+3D_0}
\end{equation}
has positive valuation. Thus $e$ is a square. If $v_p(\tau/h)\ge1$,
including $\tau=0$, the identity
\[
 -3e-1=\frac{-4(\tau/h)^2}{(\tau/h)^2+3D_0}
\]
shows that $-3e$ is a square. In either case $A$ or $-3A$ is a square,
giving a point with $q=0$.

It remains that $v_p(D)$ is even and neither $D$ nor $-3D$ is a
square. Remove an even power of $p$ from $D$, scaling $\tau$ in the
same way, to make $D$ a unit. Both $D$ and $-3D$ then have nonsquare
residues. For an integral scaled $\tau$, neither numerator nor denominator
of $e$ can vanish modulo $p$; for a nonintegral scaled $\tau$ they have
the same valuation. Therefore $v_p(e)=0$.

We claim that $v_p(C)$ is even. If it is negative, it equals the negative
valuation of $b$, which is even. If it were positive, then $b$ would be a
unit congruent to $17$ modulo $p$. The number $17b$ would be a square in
$\Q_p$, and so would $D=17b/t^2$, a contradiction. The remaining
possibility is $v_p(C)=0$. Thus $v_p(A)$ is even. Remove this even power
by scaling $s$ and $q$. The resulting quartic has unit leading coefficient
and unit $A$, so \eqref{eq:disc} is a unit. Its projective reduction is a
smooth genus-one curve over $\F_p$. It has an $\F_p$-point
\cite[Chapter~IV, Corollary~9.3]{Milne}; equivalently its number of points
satisfies the genus-one Hasse--Weil bound. Smooth lifting gives a
$\Q_p$-point. This argument includes primes dividing the original conic
denominator: if its valuation changes, one of the earlier square cases
applies.

\subsection{The primes 2, 3, 17 and the real place}

At $2$, the numbers $d$ and $17$ are squares in $\Q_2$, so the points
at infinity are rational. At the real place the same conclusion follows
from $D>0$.

At $3$, choose $h\in\Q_3^\times$ with $h^2=d$. Replacing $q$ by
$q/h$ and writing $u=\tau/h$ reduces the leading coefficient to $17$
and gives
\[
                     e=\frac{u^2-17}{u^2+51}.
\]
Preparation gives $v_3(b)\ge2$, because this valuation is even. Thus
$C=17^2(1-b/17)$ is a square unit and is $1$ modulo $3$.
If $v_3(u)<0$, then $e\equiv1\pmod3$ is a square unit.
If $v_3(u)>0$, including $u=0$, then $-3e\equiv1\pmod3$ is a
square unit. These two cases give $q=0$ as before. If $u$ is a unit,
then $e\equiv2\pmod3$, and at $s=1$ the right side of the scaled
equation is
\[
                  17(1-eC)(1+3eC)\equiv1\pmod3.
\]
It is a nonzero square in $\Q_3$.

At $17$, choose $h^2=d$ in $\Q_{17}$ and make the same scaling.
Now $v_{17}(17)=1$, and $C=17^2(1-b/17)$ is a square because
$v_{17}(b)\ge2$. The odd-valuation argument above, with $D_0=17$,
gives a point for every value of $u=\tau/h$.

This proves local solubility at every place. Lemma~\ref{lem:open} gives
the affine, nonzero-ordinate and avoidance assertions, and completes the
proof of Theorem~\ref{thm:local}.

\section{Finite families and conjugate guards}

\begin{corollary}\label{cor:guards}
Let $B_0$ be a finite etale $\Q$-algebra with a rational character
$\chi:B_0\to\Q$. For $1\le i\le r$, let $b_i,d_i,t_i$ be units with
$b_i=d_it_i^2$, and suppose each $\chi(b_i)$ is prepared. There is a
finite set $\mathcal E\subset\Q$, not containing zero, such that for
every $\tau\in\Q\setminus\mathcal E$ and every completion $k$ of
$\Q$, the $r$ pole equations at $\chi$ have affine $k$-solutions with
nonzero ordinates and
\[
 \Nm_{B_0/\Q}(\tau^2+51d_i)\ne0,
 \qquad
 \Nm_{B_0/\Q}\!\left(
 17d_i(s_i^2-A_i)(s_i^2+3A_i)\right)\ne0.
\]
Here $A_i=17(17-b_i)(\tau^2-17d_i)/(\tau^2+51d_i)$.
Every native factor guard obtained by syntactic square extraction of these
expressions is nonzero as well.
\end{corollary}
\begin{proof}
The norm of each conic denominator is a monic nonzero polynomial in
$\tau$, and its value at zero is $\Nm_{B_0/\Q}(51d_i)\ne0$.
Let $\mathcal E$ consist of their rational zeros. Fix a permitted $\tau$.
Each denominator is a unit in $B_0$. For each $i$, the radicand in the
statement is a degree-four polynomial in $s_i$ with unit leading
coefficient $17d_i$. Its norm is therefore a nonzero polynomial in
$s_i$, of degree $4\dim_{\Q}B_0$.

Apply Theorem~\ref{thm:local} to the rational numbers
$\chi(b_i),\chi(d_i),\chi(t_i)$ and to the same fixed $\tau$.
For each $i$, avoid the finitely many zeros of its norm polynomial.
These choices are independent once $\tau$ is fixed. A nonzero product
has no zero factor, which also proves the assertions for all factors
extracted from the numerator, denominator and square multiplier.
\end{proof}

In particular $\tau=0$ always supplies simultaneous local witnesses with
these conjugate guards. The corollary concerns the pole stage: it assumes
that all earlier point-algebra guards have already been met. It supplies
no simultaneous rational witness.

\section{The exact global condition}

For $n\in\Q^\times$, write
\[
                        E_n:\quad y^2=x(x-n)(x+3n).
\]
Put $n=DA_\tau$. The polynomial map
\begin{equation}\label{eq:ellmap}
                    (s,q)\longmapsto (Ds^2,Dsq)
\end{equation}
sends $\mathcal C_\tau$ to $E_n$ and extends to a degree-two map on
the smooth projective curves.

\begin{samepage}
\begin{proposition}\label{prop:global}
For prepared $b$, the following are equivalent for a fixed $\tau\in\Q$:
\begin{enumerate}
\item The pole quartic has an affine rational point with $q\ne0$.
\item $E_{DA_\tau}(\Q)$ has a point with nonzero ordinate and abscissa
in $D\Q^{\times2}$.
\item $E_{N(\tau)}(\Q)$ has such a point, where
\[
             N(\tau)=DC(\tau^2-D)(\tau^2+3D).
\]
\end{enumerate}
\end{proposition}
\end{samepage}
\begin{proof}
If $s=0$, the right side of \eqref{eq:quartic} is negative over
$\mathbb R$, because $D>0$. Hence a rational affine point has $s\ne0$,
and \eqref{eq:ellmap} proves the first implication. Conversely, write
$x=Ds^2$ with $s\in\Q^\times$ and set $q=y/(Ds)$. Substitution gives
\eqref{eq:quartic} and preserves nonzero ordinates. Finally,
\[
        N(\tau)=DA_\tau(\tau^2+3D)^2.
\]
Scaling elliptic coordinates by $(\tau^2+3D)^2$ and
$(\tau^2+3D)^3$ identifies the two curves and preserves abscissa
square classes.
\end{proof}

Thus simultaneous rational realization asks for one rational $\tau$ at
which specified abscissa classes occur on all the correlated quartic-twist
families $E_{N_i(\tau)}$. The local theorem verifies all local conditions
for each of these classes, for every rational $\tau$. It does not establish
their simultaneous rational occurrence.

This is consistent with the fixed-parameter examples of \cite{Fixed}.
Those exclude rational points by exact elliptic rank and square-class
calculations. The conic coordinates in the two normalizations satisfy
$\tau_{\rm old}=t\tau$ \cite[Section~4]{UnitNorm}, so
Theorem~\ref{thm:local} also applies to those fixed old parameters.
The associated quartics are everywhere locally soluble despite their
global obstruction. Consequently local solubility of every shared fiber
does not supply the missing simultaneous rational realization.

\section{Exact verification}

The verifier
\nolinkurl{computation/pole_parameters/shared_parameter_local_v41/check.py}
checks the two unit identities, the discriminant \eqref{eq:disc}, the map
\eqref{eq:ellmap}, and the polynomial twist identity exactly. It counts
all $7704$ smooth unit quartics of the displayed form over the twelve
prime fields from $\F_5$ through $\F_{43}$ in its specified list, checks
the exceptional residue calculation at $3$, and exercises $19008$ exact
local valuation branches, including rescaled square representatives.
It also reconstructs all four current source circuits and checks their
v40 guard specifications: the sole pole-stage guard factor involving
$\tau_i$ but not $s_i$ is the conic denominator, whose value at
$\tau_i=0$ is $51$ times the numerator of the unit $d_i$.
The reconstructed first ordinate is neither stored nor guarded; the
nonzero guards on the second ordinates can be met by the local witnesses
of Theorem~\ref{thm:local}.

It also gives a local certificate at every possible bad prime for the
published prepared inputs excluding old $\tau=\pm1$ and $\tau=\pm17$,
for the $e=-1/3$ counterexample, and for an input of square class $409$.
At all remaining primes the corresponding quartic has smooth genus-one
reduction. An invalid-preparation control retains all exceptional local
square conditions and valuation inequalities but replaces $b=(204/55)^2$
by $b/409$. Its negative odd valuation at $409$ makes the quartic at
$\tau=0$ locally insoluble there, by an exact valuation-parity argument.
These computations verify explicit identities and instances;
the proof in Section~2 establishes the assertion for every prepared input
and every rational parameter. No rank computation or local-to-global
assumption enters the local theorem.

\begin{thebibliography}{9}
\bibitem{Fixed}
\emph{Obstructions to fixed parameters in pole certificates},
\nolinkurl{papers/pole_parameters/fixed_pole_parameters_20261109_v17.tex}.
\bibitem{Milne}
J. S. Milne, \emph{Elliptic Curves}, second edition,
World Scientific, 2021. Chapter~IV, Corollary~9.3 and Theorem~9.4.
\url{https://jmilne.org/math/Books/EC2.pdf}.
\bibitem{UnitNorm}
\emph{A single norm determinant for the 25- and 13-variable constructions},
\nolinkurl{papers/finite_algebra/unit_norm_20261114_v28.tex}.
\end{thebibliography}
\end{document}
